Though our work has application beyond scRNA-seq, we next describe our data and goals using scRNA-seq terminology to clarify and concretize our notation.

\textbf{Data Setup.}  
We consider single-cell mRNA measurements collected at $\totsteps$ distinct time points, labeled $\timestep{1} < \timestep{2} < \cdots < \timestep{\totsteps}$. For convenience, we set $\timestep{1} = 0$. We do not require these time points to be equally spaced. At each time $\timestep{\timeidx}$, the observed data consist of $\tottrajec_{\timeidx}$ cells, each providing a single mRNA expression level measurement (representing a cell state) in $\mathbb{R}^d$, denoted $\obsat{\timestep{\timeidx}}^{\sampleidx}$. After a cell's mRNA level is measured, the cell is destroyed. So each cell appears exactly once in the dataset. We therefore collect $\tottrajec = \sum_{\timeidx=1}^{\totsteps} \tottrajec_{\timeidx}$ total observations across all time points. We write $\obsatall{\timestep{\timeidx}}$ for the full set of measurements taken at time $\timestep{\timeidx}$.

\textbf{Goal.} If a cell's mRNA expression level were not measured, the cell would remain alive, and its mRNA expression would evolve continuously over time along a (latent) trajectory. 
Formally, we denote the latent trajectory of the $\sampleidx$th cell observed at the $\timeidx$th time step by $\responseat{\timet}^{(\timeidx, \sampleidx)}$. The observed state of the cell at time $\timestep{\timeidx}$ is a single point on this trajectory $\obsat{\timestep{\timeidx}}^{\sampleidx} = \responseat{\timet = \timestep{\timeidx}}^{ (\timeidx, \sampleidx) }$. We assume these latent trajectories are independent realizations from an underlying latent distribution. 
Consequently, the observed snapshot measurements are also independent.
Our objective is
to infer a probabilistic model, chosen from a specified parametric family, that best captures the distribution of these unobserved trajectories. Our model should provide a distribution over forecasted trajectories beyond the measured time and also over interpolated trajectories between observed times. 


%\textbf{Dataset and independence assumptions.} By the assumptions above, the dataset $\{\obsat{\timestep{\timeidx}}^{\sampleidx}\}_{\timeidx,\sampleidx}$ represents independent observations. But cells measured at different times can be expected to reflect different state distributions. So, if we treat the observation times of data points as fixed and known, it's not reasonable to assume the observations are both independent and identically distributed (i.i.d.). However, a key insight is that we can treat the observation times 
%themselves as random draws from an underlying distribution $\timemeasure(\timestep{})$. Then a dataset of pairs $\{(\obsat{\timestep{\timeidx}}^{\sampleidx}, \timestep{\timeidx})\}_{\timeidx,\sampleidx}$ can be treated as i.i.d.\ samples from a joint state-time distribution. While it follows that $\{\obsat{\timestep{\timeidx}}^{\sampleidx}\}_{\timeidx,\sampleidx}$  are now i.i.d.\ as well, using the pairs $\{(\obsat{\timestep{\timeidx}}^{\sampleidx}, \timestep{\timeidx})\}_{\timeidx,\sampleidx}$ as our data aligns better with real biological experiments. Instead of sequencing each sample immediately after collection, experimenters often tag each sample with a unique identifier that encodes the time it was collected. Then, all cells are sequenced together in a single batch, with the time information retrieved from the tags. 

\textbf{Dataset and independence assumptions.} 
At each observation time $\timestep{\timeidx}$, the measurements 
$\{\obsat{\timestep{\timeidx}}^{\sampleidx}\}_{\sampleidx}$ are modeled as i.i.d.\ draws from a conditional distribution $p(Y \mid t = \timestep{\timeidx})$. 
Across different timepoints, say $\timestep{\timeidx} \neq \timestep{j}$, the conditional distributions $p(Y \mid t = \timestep{\timeidx})$ and $p(Y \mid t = \timestep{j})$ may differ, so the overall dataset is independent but not identically distributed if we treat times as fixed. 
But we can model the observation times themselves as random draws from a discrete distribution $\timemeasure(t)$ supported on the set of measurement times. 
Then the dataset of pairs $\{(\obsat{\timestep{\timeidx}}^{\sampleidx}, \timestep{\timeidx})\}_{\timeidx,\sampleidx}$ can be regarded as i.i.d.\ samples from the joint distribution $p(Y, t) = p(Y \mid t)\timemeasure(t)$. 
This joint view aligns with real biological experimental protocols. Indeed, instead of sequencing each sample immediately after collection, experimenters often tag each sample with a unique identifier that encodes the time it was collected. Then, all cells are sequenced together in a single batch, with the time information retrieved from the tags. In \cref{app:iid-assumptions} we discuss situations where this i.i.d.\ assumption might not hold.


\textbf{Model.} We model each latent trajectory of the $(\timeidx, \sampleidx)$ cell with a stochastic differential equation (SDE) driven by a $d$-dimensional Brownian motion $\brownianm^{ (\timeidx,\sampleidx) }$, independent across particles:
\begin{equation}
    \de\responseat{\timet}^{ (\timeidx,\sampleidx) } = \truedrift(\responseat{\timet}^{ (\timeidx,\sampleidx) },t) \de\timet + \truevolatility(\responseat{\timet}^{ (\timeidx,\sampleidx) },\timet) \de\brownianm,~~\responseat{\timet = 0}^{ (\timeidx,\sampleidx) }\sim \marginal{0}.
    \label{eq:mainsde}
\end{equation}
We assume that the drift $\truedrift(\cdot,\cdot): \R^{d}\times [0,\timet_{\totsteps}] \to \R^{d}$ and initial marginal distribution $\marginal{0}$ are unknown. 
Previous Schrödinger bridge methods typically assume fixed, known volatility \citep[e.g.][]{de2021diffusion,Vargas2021,koshizuka2022neural,Wang2023,shen2024multi,zhang2024joint,guan2024identifying,chen2024deep}. By contrast, we allow the common case where the volatility function $\truevolatility$ can be unknown and also state- and time-dependent.
For example, in scRNA-seq, transcriptional noise varies with gene identity, cell state, and developmental stage --- and is further confounded by technical artifacts such as amplification bias and stochastic capture. Volatility in these systems reflects true biological uncertainty and is rarely known in advance. 

Finally, we assume standard regularity conditions on the SDE: Lipschitz continuity and linear growth for drift and volatility \citep[to ensure existence of strong solutions to the SDE; see][Chapter 3, Theorem 3.1]{pavliotis2016stochastic} and bounded second moments of the particle distributions (ruling out the possibility that the process exhibits unbounded variability). We state these assumptions formally in \cref{app:assumptions}.

\textbf{One natural extension of Schrödinger bridges to forecasting highlights limitations of fixed reference dynamics.} 
SB methods reconstruct distributions of trajectories by matching observed marginals while penalizing deviation from a predefined reference process, typically Brownian motion; see \cref{eq:sb-problem} in \cref{app:sb-forecasting} for a standard setup. These methods are well suited for interpolation between observed marginals. Forecasting, however, has not been a  focus in this line of work. In \cref{app:sb-forecasting}, we work through one natural extension of the SB framework to forecasting by introducing an additional, future marginal. In this setup, under fixed reference dynamics, extrapolation reduces to propagating the final observed marginal forward under the reference process. But the reference is chosen without considering any data, so we expect poor extrapolation performance.
%SB methods reconstruct distributions of trajectories by solving a constrained optimization problem that matches observed marginals while penalizing deviation from a predefined reference process, typically Brownian motion; see \cref{eq:sb-problem} in \cref{app:sb-forecasting} for a detailed formulation. While these methods are well suited for interpolation between observed marginals, their forecasting ability is more limited: in many formulations, extrapolation amounts to propagating the final observed marginal forward under the reference dynamics. Extensions such as \citet{chen2024deep} introduce velocity or momentum to improve extrapolation, but they still rely on fixed reference dynamics and do not learn nonlinear, state-dependent dynamics from data. As a result, their predictive power remains constrained by the choice of reference or by restrictive optimization objectives. 
